\documentclass[12pt]{article}
\usepackage[utf8]{inputenc}
\usepackage[spanish]{babel}
\usepackage{amsmath, amssymb, amsthm}
\usepackage{tcolorbox}
\usepackage{geometry}
\geometry{margin=2.5cm}

\title{Unidad 6: Anillos y aritmética modular}
\author{Matemática Discreta - UADER 2025}
\date{}

\begin{document}

\maketitle

\section*{Estructura de un anillo}

\begin{tcolorbox}[title=Definición de anillo]
Sea \( R \) un conjunto no vacío con dos operaciones binarias, comúnmente denotadas por \( + \) y \( \cdot \), que pueden ser diferentes de las operaciones usuales. El conjunto \( (R, +, \cdot) \) es un \textbf{anillo} si para todos \( a, b, c \in R \) se cumplen las siguientes condiciones:
\begin{itemize}
    \item \( a + b = b + a \) (\textbf{conmutativa} de la suma)
    \item \( a + (b + c) = (a + b) + c \) (\textbf{asociativa} de la suma)
    \item Existe un elemento \( 0 \in R \) tal que \( a + 0 = 0 + a = a \) (\textbf{identidad aditiva})
    \item Para cada \( a \in R \), existe \( -a \in R \) tal que \( a + (-a) = (-a) + a = 0 \) (\textbf{inverso aditivo})
    \item La multiplicación es \textbf{asociativa}: \( a \cdot (b \cdot c) = (a \cdot b) \cdot c \)
    \item La multiplicación distributiva respecto a la suma:
    \[
    a \cdot (b + c) = a \cdot b + a \cdot c, \quad (a + b) \cdot c = a \cdot c + b \cdot c
    \]
\end{itemize}
\end{tcolorbox}

\begin{tcolorbox}[colback=yellow!15, title=Notas y aclaraciones]
\begin{itemize}
    \item Cuando se trabaja con la segunda operación, es común denotar el producto como \( ab \).
    \item Por inducción matemática, estas propiedades se extienden para cualquier número finito de elementos.
    \item \textbf{Nota personal:} Es fundamental verificar la cerradura bajo ambas operaciones, especialmente en subconjuntos propuestos como subanillos o ideales.
\end{itemize}
\end{tcolorbox}

\section*{Propiedades adicionales}

\begin{tcolorbox}[title=Propiedades importantes]
En cualquier anillo \( (R, +, \cdot) \):
\begin{itemize}
    \item La identidad aditiva, si existe, es única.
    \item El inverso aditivo de cada elemento también es único.
    \item La identidad multiplicativa, si existe, es única.
    \item En un anillo con unidad, un elemento \( u \) es \textbf{unidad} o \textbf{invertible} si existen \( a \in R \) tal que \( au = ua = 1 \).
    \item Un \textbf{dominio de integridad} es un anillo conmutativo, con unidad, y sin divisores propios de cero.
    \item Un \textbf{cuerpo} (o campo) es un anillo con unidad en el que todo elemento diferente de cero tiene inverso multiplicativo.
\end{itemize}
\end{tcolorbox}

\section*{Subanillos e ideales}

\begin{tcolorbox}[title=Subanillo]
Sea \( (R, +, \cdot) \) un anillo. Un subconjunto no vacío \( S \subseteq R \) es un \textbf{subanillo} si:
\begin{itemize}
    \item \( S \) es cerrado respecto a la suma: \( a + b \in S \) para \( a, b \in S \).
    \item \( S \) es cerrado respecto a la multiplicación: \( a b \in S \) para \( a, b \in S \).
\end{itemize}
\end{tcolorbox}

\begin{tcolorbox}[title=Ideales]
Un subconjunto no vacío \( I \subseteq R \) es un \textbf{ideal} si:
\begin{itemize}
    \item \( a - b \in I \) para \( a, b \in I \).
    \item Para todo \( r \in R \), se verifica que \( r a \in I \) y \( a r \in I \).
\end{itemize}
\end{tcolorbox}

\begin{tcolorbox}[colback=yellow!15, title=Nota personal]
\textbf{Importante:} Todo ideal es un subanillo, pero no todos los subanillos son ideales. La diferencia radica en la cerradura multiplicativa del lado del elemento externo \( r \).
\end{tcolorbox}

\section*{Los enteros módulo \( n \)}

\begin{tcolorbox}[title=Congruencia modular]
Sea \( n \in \mathbb{N} \), \( n > 1 \). Para \( a, b \in \mathbb{Z} \), decimos que
\[
a \equiv b \pmod{n}
\]
si \( n \mid (a - b) \), es decir, si \( a - b = kn \) para algún \( k \in \mathbb{Z} \). 
\textbf{Nota:} esta relación de equivalencia divide a \( \mathbb{Z} \) en \( n \) clases de equivalencia, conocidas como residuos o restos módulo \( n \).
\end{tcolorbox}

\begin{tcolorbox}[colback=yellow!15, title=Nota personal]
En el estudio de los residuos módulo \( n \), se trabaja con los representantes en \( \{0, 1, ..., n-1\} \). Es crucial recordar que las operaciones están bien definidas en estas clases de equivalencia, formando un anillo \( \mathbb{Z}_n \).
\end{tcolorbox}

\section*{Operaciones en \( \mathbb{Z}_n \)}

\begin{align*}
[a]_n + [b]_n &= [a + b]_n \\
[a]_n \cdot [b]_n &= [a b]_n
\end{align*}

Estas definiciones hacen que \( \mathbb{Z}_n \) sea un anillo con \( n \) elementos. Además:
\begin{itemize}
    \item \( \mathbb{Z}_p \) es un cuerpo si y solo si \( p \) es primo.
    \item Los elementos \( a \in \mathbb{Z}_n \) que son coprimos con \( n \) son unidades.
    \item El número total de unidades en \( \mathbb{Z}_n \) es \( \varphi(n) \), donde \( \varphi \) es la función de Euler.
\end{itemize}

\begin{tcolorbox}[title=Notas y aplicaciones]
La aritmética modular es fundamental en criptografía, teoría de códigos, y en la resolución de problemas en álgebra abstracta. Es esencial manejar con precisión las congruencias y su relación con los ideales en anillos cocientes.
\end{tcolorbox}

\section*{Homomorfismos e isomorfismos}

\begin{tcolorbox}[title=Definición]
Sean \( (R, +, \cdot) \) y \( (S, \oplus, \otimes) \) anillos. Una función \( f: R \to S \) es un \textbf{homomorfismo de anillos} si:
\[
f(a + b) = f(a) \oplus f(b), \quad f(a b) = f(a) \otimes f(b)
\]
Si \( f \) es biyectiva, entonces es un \textbf{isomorfismo}.
\end{tcolorbox}

\section*{El Teorema del Resto Chino}

\begin{tcolorbox}[title=Teorema]
Sea \( n = p_1^{e_1} p_2^{e_2} \cdots p_k^{e_k} \) la factorización en primos de \( n \). Entonces:
\[
\mathbb{Z}_n \cong \mathbb{Z}_{p_1^{e_1}} \times \mathbb{Z}_{p_2^{e_2}} \times \cdots \times \mathbb{Z}_{p_k^{e_k}}
\]
y las operaciones se definen componente a componente. Además, los anillos \( \mathbb{Z}_{n} \) y los productos anteriores son ísomorfos.
\end{tcolorbox}

\begin{tcolorbox}[colback=yellow!15, title=Notas y aplicaciones]
Este resultado permite descomponer problemas en módulos primos y analizar su comportamiento en estructuras más simples. Es muy útil en criptografía y algebra computacional.
\end{tcolorbox}

\end{document}
